\subsection{Suposiciones y Dependencias}{\color{black}
Se supone que los datos territoriales, meteorológicos y de vegetación disponibles para el sector de Aguas Claras representan de manera adecuada las condiciones que se incorporarán a las simulaciones. Entre estos se encuentran la información del terreno y las características ambientales empleadas para configurar los escenarios, cuya calidad y cobertura condicionan la fidelidad de los resultados. Para la consulta de una emergencia, se supone que las condiciones meteorológicas obtenidas mediante Open-Meteo o ingresadas por el usuario son correctas y corresponden a la situación que se desea analizar. La aplicación depende de la disponibilidad del servicio cuando se utilice la consulta automática, aunque permite el ingreso manual si no hay conexión o el servicio no está disponible.

La construcción \textcolor{black}{del dataset} depende de la disponibilidad de Fire Dynamics Simulator (FDS) y del acceso a infraestructura de cómputo de alto desempeño (HPC-ZINE) con capacidad para ejecutar las simulaciones previstas. Se supone que estos recursos estarán disponibles durante el desarrollo, puesto que los escenarios deben calcularse previamente y almacenarse junto con sus condiciones de entrada para que la aplicación pueda consultarlos~\cite{vfp2026}. Si los datos requeridos o la infraestructura no están disponibles, la cantidad y cobertura de los escenarios que podrán incorporarse \textcolor{black}{al dataset} se verán condicionadas.
}
